\PassOptionsToPackage{table,xcdraw}{xcolor}
\documentclass[10pt, a4paper, landscape]{article}

\usepackage[T1]{fontenc}
\usepackage[utf8]{inputenc}
\usepackage{tgpagella}
\usepackage[spanish, es-noshorthands]{babel}
\usepackage{amsmath, amssymb}
\usepackage{booktabs}
\usepackage{tabularx}
\usepackage{tikz}
\usepackage[american]{circuitikz}
\usetikzlibrary{arrows.meta, positioning, calc}
\usepackage[margin=1.5cm]{geometry}
\usepackage{enumitem}
\usepackage{microtype}

\emergencystretch=1.5em
\definecolor{ucbblue}{RGB}{0, 51, 102}
\definecolor{ucbgray}{RGB}{240, 244, 248}
\definecolor{warnred}{RGB}{180, 20, 20}

\begin{document}

\noindent
{\LARGE \textbf{\textcolor{ucbblue}{IMT-221 — Circuitos Electrónicos III}}}\\[0.12cm]
{\Large \textbf{Hito 3: Guía Única E1 $\rightarrow$ E4 — Parte 2: Instrumentación y Etapa E1 (NMOS Low-Side Resistivo)}}\\[0.08cm]
{\small \textbf{Versión integrada 1.3} $\cdot$ 6 de noviembre de 2026 $\cdot$ Condicionado a firma de Puerta G1}
\par\vspace{0.15cm}\noindent\rule{\linewidth}{0.4pt}\vspace{0.2cm}

\section*{2. Esquema Común e Instrumentación}

\subsection*{2.1 Circuito E1 (Mantenido también como carga en E2)}

\begin{center}
\begin{circuitikz}[scale=0.92, transform shape, thick]
    % Alimentación y carga resistiva
    \draw (0,4.4) node[above]{\textbf{$+12\,\text{V}_{\text{BUS}}$} (Fuente limitada)} to[R, l={$R_L = 100\,\Omega$}, a={\small Potencia verificada}, i={$I_D$}] (0,2.0);

    % Transistor NMOS M1
    \draw (0,0.2) node[nigfete] (M1) {};
    \node[above left=0.05cm] at (M1.G) {G};
    \node[above right=0.05cm] at (M1.D) {D};
    \node[below right=0.05cm] at (M1.S) {S};
    \node[below right=0.05cm] at (M1.B) {B};

    \coordinate (drain) at ($(M1.D) + (0, 0.55)$);
    \coordinate (source) at ($(M1.S) + (0, -0.75)$);
    \coordinate (gate) at ($(M1.G) + (-2.2, 0)$);
    \coordinate (ctrl) at ($(gate) + (-3.0, 0)$);
    \coordinate (gnd_ctrl) at (ctrl |- source);
    \coordinate (gnd_bus) at ($(gnd_ctrl) + (0, -0.75)$);

    % Conexiones Drain, Source y Bulk
    \draw (0,2.0) -- (drain) node[circle, fill, inner sep=1.5pt]{} -- (M1.D);
    \draw (M1.S) -- (source) node[ground]{};
    \draw (M1.B) -- ++(0.6,0) |- (source);

    % Diodo intrínseco explícito a la derecha
    \draw[blue, dashed] (M1.S) -- ++(2.2,0) coordinate (db_s) to[D, l_={\small Diodo intrínseco}] (db_s |- M1.D) -- (M1.D);

    % Sondas de osciloscopio
    \draw[-{Stealth[length=2.5mm]}, warnred, very thick] (drain) -- ++(3.2,0) node[right]{\textbf{CH2}: Punta en D, ref.\ en S ($V_{DS}$)};
    \draw[-{Stealth[length=2.5mm]}, warnred, very thick] (gate) -- ++(0,1.5) node[above]{\textbf{CH1}: Punta en G, ref.\ en S ($V_{GS}$)};

    % Red de excitación de compuerta
    \draw (M1.G) -- (gate) node[circle, fill, inner sep=1.5pt]{};
    \draw (gate) to[R, l_={$R_G = 220\,\Omega$}] (ctrl) node[circle, fill, inner sep=1.5pt]{} node[above]{\textbf{Control}};
    \draw (ctrl) to[V, l_={\small $0\text{--}5\,\text{V}$ / $0\text{--}10\,\text{V}$}] (gnd_ctrl) node[ground]{};
    \draw (gate) to[R, l={$R_{PD} = 100\,\text{k}\Omega$}] (gate |- source) node[ground]{};

    % Retorno común verificado
    \draw[dashed, ucbblue] (gnd_bus) -- (source |- gnd_bus) node[midway, below]{\small Retorno común GND potencia / control (unir solo tras revisar tierras)};
\end{circuitikz}
\end{center}

\noindent
\textbf{Reglas de instrumentación física:}
\begin{itemize}[noitemsep, topsep=2pt]
    \item \textbf{Revisar físicamente G/D/S y pestaña TO-220 antes de cablear.} Identificar qué tierra está conectada a protección; un osciloscopio de banco suele tener las masas de las sondas conectadas a tierra.
    \item \textbf{CH1:} punta gate, referencia source, mide $V_{GS}$ en low-side. \textbf{CH2:} punta drain, referencia source, mide $V_{DS}$. \textbf{\textcolor{warnred}{No conectar la pinza de tierra al drain ni al lado superior de la carga.}}
    \item Medir la corriente mediante instrumento/método aprobado; un multímetro en modo amperímetro va \textbf{en serie}, no en paralelo. Registrar escala, atenuación, acoplamiento y límites.
\end{itemize}

\subsection*{2.2 Secuencia Universal Antes de una Energización (Todas las Etapas)}
\begin{enumerate}[noitemsep, topsep=2pt]
    \item Dibujar circuito y comprobar continuidad con \textbf{todas las fuentes apagadas}.
    \item Identificar componentes reales, pinout y límite de cada fuente, conductor y sonda.
    \item Preparar LTspice: guardar esquema, modelo/fuente, tiempo de transitorio y log.
    \item Verificar \textbf{señal de control aislada de la potencia} en osciloscopio (mínimo, máximo, frecuencia, duty y flancos).
    \item Unir tierras \textbf{solo después} de revisar si hay conexión a tierra física y la topología.
    \item Comprobar que el gate tiene estado OFF definido al desconectar señal, y que no se excede $V_{GS(\max)}$.
    \item Presentar la puerta G correspondiente al docente; anotar \texttt{autorizado / observado / bloqueado}, fecha y firma.
    \item \textbf{Solo si autorizado:} limitar corriente de fuente, comenzar en estado OFF, avanzar por estados aprobados, medir, detener ante calor/corriente/tensión fuera de rango.
    \item Apagar alimentación y descargar nodos antes de modificar cables. Guardar evidencia y discrepancias.
\end{enumerate}

\newpage
\section*{3. Etapa E1 — Interruptor NMOS Low-Side Resistivo}

\subsection*{3.1 Predicción Matemática Mínima (Derivación Paso a Paso)}

\paragraph{1. Estado ON Ideal ($R_{DS(\text{on})} \rightarrow 0\,\Omega$, $V_{DS,\text{ON}}^{\text{ideal}} = 0\,\text{V}$):}
Para el escenario de estudio $V_{\text{BUS}} = 12\,\text{V}$ y $R_L = 100\,\Omega$, aplicando la Ley de Tensiones de Kirchhoff (KVL) al lazo principal de potencia:
\begin{align}
    V_{\text{BUS}} - I_{D,\text{ON}}^{\text{ideal}} R_L - V_{DS,\text{ON}}^{\text{ideal}} &= 0 \implies I_{D,\text{ON}}^{\text{ideal}} = \frac{V_{\text{BUS}}}{R_L} = \frac{12\,\text{V}}{100\,\Omega} = 0.12\,\text{A} = 120\,\text{mA}, \\
    P_{RL,\text{ON}}^{\text{ideal}} &= \left(I_{D,\text{ON}}^{\text{ideal}}\right)^2 R_L = (0.12\,\text{A})^2 \times 100\,\Omega = 0.0144\,\text{A}^2 \times 100\,\Omega = 1.44\,\text{W}.
\end{align}

\paragraph{2. Estado OFF Ideal y Régimen PWM ($f = 100\,\text{Hz}$, duty $D = 0.50$):}
En OFF ideal, $I_D \approx 0\,\text{A}$ y $V_{DS} \approx 12\,\text{V}$; en ON ideal, $V_{DS} \approx 0\,\text{V}$. A $f = 100\,\text{Hz}$, el periodo es $T = \frac{1}{100\,\text{Hz}} = 10\,\text{ms}$, con $t_{\text{ON}} \approx t_{\text{OFF}} \approx 5\,\text{ms}$. La potencia \textbf{media ideal} de la carga es:
\begin{equation}
    P_{RL,\text{med}}^{\text{ideal}} = \frac{1}{T}\int_0^T p_{RL}(t)\,dt = D \times P_{RL,\text{ON}}^{\text{ideal}} = 0.50 \times 1.44\,\text{W} = 0.72\,\text{W}.
\end{equation}
La resistencia debe soportar también transitorios, temperatura y posibles estados ON continuos ($1.44\,\text{W}$); \textbf{$1/4\,\text{W}$ ($0.25\,\text{W}$) es inadecuada}.

\paragraph{3. Estado ON No Ideal (con $R_{DS(\text{on})}$ aplicable a ese $V_{GS}$ y temperatura):}
\begin{align}
    I_D &\approx \frac{V_{\text{BUS}}}{R_L + R_{DS(\text{on})}}, \qquad V_{DS,\text{on}} \approx I_D R_{DS(\text{on})}, \qquad P_{\text{MOS,cond}} \approx I_D^2 R_{DS(\text{on})}.
\end{align}
Para $R_{DS(\text{on})} = 25\,\text{m}\Omega = 0.025\,\Omega$ \textbf{hipotético y condicionado}:
\begin{align*}
    I_D &\approx \frac{12\,\text{V}}{100\,\Omega + 0.025\,\Omega} = \frac{12\,\text{V}}{100.025\,\Omega} \approx 0.119970\,\text{A} = 119.97\,\text{mA}, \\
    V_{DS,\text{on}} &\approx 0.119970\,\text{A} \times 0.025\,\Omega \approx 0.002999\,\text{V} \approx 3.00\,\text{mV}, \\
    P_{\text{MOS,cond}} &\approx (0.119970\,\text{A})^2 \times 0.025\,\Omega \approx 0.0003598\,\text{W} \approx 0.36\,\text{mW}.
\end{align*}
\textbf{No presentar $3\,\text{mV}$ como lectura garantizada:} puede quedar debajo de la incertidumbre de los instrumentos de aula. \textbf{IRLZ44N:} el fabricante Infineon especifica conducción a niveles lógicos (el sitio oficial declara $R_{DS(\text{on})}$ máx.\ a $4.5\,\text{V}$). \textbf{IRFZ44N:} el $R_{DS(\text{on})}$ máximo de $17.5\,\text{m}\Omega$ se publica a $V_{GS}=10\,\text{V}$; con gate a $5\,\text{V}$ \textbf{puede conducir} esta carga, pero esa ficha no garantiza la misma resistencia. $V_{GS(\text{th})}$ tampoco significa transistor completamente encendido.

\subsection*{3.2 Señal de Control (Misma Frecuencia para Todos)}
\begin{itemize}[noitemsep, topsep=2pt]
    \item \textbf{Generador principal:} $100\,\text{Hz}$, cuadrada, $50\%$. Para \texttt{IRLZ44N} buscar \textbf{$0\text{--}5\,\text{V}$}; para \texttt{IRFZ44N} buscar \textbf{$0\text{--}10\,\text{V}$}. Si se dispone de offset, señal de $10\,\text{V}_{\text{pp}}$ y $+5\,\text{V}$ offset es el objetivo matemático de $0/10\,\text{V}$, \textbf{sujeto a la escala y carga del equipo}. Medir siempre sin gate y luego con gate. Si un \texttt{SKMI SK3000} entrega señal bipolar $-5/+5\,\text{V}$ y no tiene offset, \textbf{no} se conecta como si fuera $0/10\,\text{V}$. Su folleto indica MAIN $50\,\Omega$ y no especifica offset ajustable en el modelo SK3000.
    \item \textbf{Arduino Uno $5\,\text{V}$ (opcional):} Control experimental de $100\,\text{Hz}$, duty $50\%$, con $R_G$ inicial \textbf{$330\,\Omega$} sujeto a verificación de corriente GPIO y flancos, $R_{PD}=100\,\text{k}\Omega$, tierra común revisada. Para \texttt{IRLZ44N} puede probarse con G1; para \texttt{IRFZ44N} es solo ensayo diagnóstico a $5\,\text{V}$. \textbf{ESP32 $3.3\,\text{V}$ no es la ruta obligatoria} para MOSFET de E1. Si no se logra una señal adecuada, no se improvisa un LM741/LM358 de gate: se consulta al docente y se usa E2 como problema de diseño.
\end{itemize}

\noindent\textbf{Código de demostración UNO ($\approx 100\,\text{Hz}$, no PWM de periférico):}
\begin{verbatim}
const byte PIN_GATE=9;
void setup(){ digitalWrite(PIN_GATE,LOW); pinMode(PIN_GATE,OUTPUT); }
void loop(){ digitalWrite(PIN_GATE,HIGH); delay(5); digitalWrite(PIN_GATE,LOW); delay(5); }
\end{verbatim}

\subsection*{3.3 LTspice E1 — Procedimiento y 3.4 Construcción y Medición E1}
\begin{enumerate}[noitemsep, topsep=2pt]
    \item \textbf{LTspice E1:} (1) Fuente DC \texttt{VBUS=12}, \texttt{RL=100}, NMOS y retorno a ground. (2) Control \texttt{PULSE(0 5 1m 1u 1u 5m 10m)} para $5\,\text{V}$ o \texttt{PULSE(0 10 1m 1u 1u 5m 10m)} para $10\,\text{V}$; \texttt{RG=220} y \texttt{RPD=100k}. (3) Emplear macromodelo comercial si se verificó; si se usa \texttt{MDEMO}, marcarlo \textbf{modelo ficticio}, no \texttt{IRLZ44N}/\texttt{IRFZ44N} (no duplicar capacitancias ficticias en macromodelo real). (4) Ejecutar \texttt{.tran 0 31m 0 1u}; graficar \texttt{V(ctrl)}, \texttt{V(gate)}, \texttt{V(drain)}, \texttt{I(RL)} a iguales escalas. (5) Registrar retardo entre $V_{GS}$ y $V_{DS}$ y estimar $p_D(t)=V_{DS}I_D$.
    \item \textbf{Construcción y medición E1:} \textbf{Paso 1:} verificar físicamente $R_L$, MOSFET, fuente y conectores; montar low-side \textbf{sin tensión}. \textbf{Paso 2:} comprobar generador solo, luego control gate con bus apagado. \textbf{Paso 3:} solicitar G1. \textbf{Paso 4:} autorizado, verificar primero OFF y luego ON con corriente limitada; después señal $100\,\text{Hz}$. \textbf{Paso 5:} registrar $V_{GS(\text{ON/OFF})}$, $V_{DS(\text{ON/OFF})}$, $I_{D(\text{ON})}$, frecuencia, duty y temperatura. \textbf{Paso 6:} confrontar con cálculo y LTspice.
\end{enumerate}

\par\vspace{0.1cm}\noindent
\renewcommand{\arraystretch}{1.22}
\begin{tabularx}{\linewidth}{|>{\raggedright\arraybackslash}X|c|c|c|c|}
\hline
\rowcolor{ucbgray}
\textbf{Estado (Tabla de Referencia E1 — Ideal, NO Experimental)} & \textbf{$V_{GS}$ Ideal} & \textbf{$V_{DS}$ Ideal} & \textbf{$I_D$ Ideal} & \textbf{$P_{RL}$ Instantánea Ideal} \\ \hline
OFF & $0\,\text{V}$ & $12\,\text{V}$ & $0\,\text{A}$ & $0\,\text{W}$ \\ \hline
ON, \texttt{IRLZ44N} & $5\,\text{V}$ & $\approx 0\,\text{V}$ & $\approx 120\,\text{mA}$ & $\approx 1.44\,\text{W}$ \\ \hline
ON, \texttt{IRFZ44N} & $10\,\text{V}$ & $\approx 0\,\text{V}$ & $\approx 120\,\text{mA}$ & $\approx 1.44\,\text{W}$ \\ \hline
\end{tabularx}

\end{document}
